\chapter*{Conclusions and Future Work}
\label{cap:conclusions}
\addcontentsline{toc}{chapter}{Conclusions and Future Work}

The main objective of this work was to design and evaluate a heterogeneous system on the
Digilent Zybo board, using the ARM processor to control the application and the FPGA to
accelerate a computationally intensive task. After completing the development process,
it can be stated that this objective has been achieved. The final system is able to send
an image from the ARM, transfer it to the programmable logic through AXI DMA, execute a
convolutional model generated with hls4ml, and retrieve the output logits in order to
obtain the predicted class.

The path toward this final version was not straightforward. The main difficulty was not
only training a CIFAR-10 model, but making that model actually fit within a Zynq Z-7010.
Many configurations that seemed reasonable in terms of parameter count became unfeasible
after implementation in Vivado, especially due to LUT usage and slice packing issues.
This part of the work made it clear that, on a small FPGA, discussing only the model size
is not enough: the way operations are synthesized and placed on the device is just as
important.

\section*{Achievement of the objectives}
\addcontentsline{toc}{section}{Achievement of the objectives}
The first part of the project made it possible to bring up the platform and understand
the complete Vivado/Vitis workflow. Bare-metal execution on the ARM was validated, the
communication between the Processing System and the Programmable Logic was configured,
and AXI-Lite and AXI DMA were used. These steps were necessary to ensure that the
artificial intelligence component was not left as an isolated IP block without real
system-level integration.

The matrix multiplication developed in Phase 3 served as an intermediate milestone. It
was a more controlled test case than the CNN, but still sufficient to verify that Vitis
HLS could generate competitive hardware and that the ARM--DMA--FPGA data path worked
correctly. In addition, the comparison with a VHDL implementation helped assess the use
of HLS from a more informed perspective. The tool does not remove the need to understand
the underlying hardware, but it does significantly reduce the effort required to reach a
functional parallel architecture.

In Phase 4, a complete CIFAR-10 model was successfully generated and integrated. The
accuracy of the model had to be balanced against the resources available on the Zybo.
The final version is not a large network, nor is it intended to compete with modern
computer vision models, but it does fulfill the purpose of this work: demonstrating AI
inference running on dedicated hardware inside a resource-constrained FPGA.

Finally, the evaluation phase made it possible to measure the benefit of the
acceleration. FPGA inference achieved an average execution time close to 2 ms, compared
with approximately 171 ms for the hls4ml-generated \texttt{C++} core running on the ARM. This
represents an approximate speedup of 85 times. Beyond the exact numerical value, the
important point is that the acceleration was measured on the real board and includes the
communication cost with the IP core.

\section*{Lessons learned}
\addcontentsline{toc}{section}{Lessons learned}
One of the clearest conclusions of the project is that HLS simplifies development, but
does not make it automatic. Directives, data types, and the reuse factor have a direct
impact on the final resource usage. On several occasions, it was necessary to go back to
the model, modify its architecture, or adjust the quantization because the design did not
fit on the FPGA. In this regard, working with hls4ml was particularly interesting: it
brings the machine learning workflow much closer to hardware design, but it still
requires careful analysis of the synthesis, utilization, and timing reports.

It was also important to distinguish between the different validation levels. At first,
when the execution on the board produced inconsistent results, it was tempting to assume
that the model itself was faulty. However, part of the problem came from Vivado/Vitis
cache issues and older versions of the IP that were still being used during integration.

Another practical lesson was the importance of keeping projects well organized. In a
workflow involving Vivado, Vitis, Vitis HLS, trained models, generated IP cores, and
auxiliary scripts, it is easy to end up with duplicated folders or old results mixed
with newer versions. Separating the project into different phases, each with its own
Vivado/Vitis project folder, makes the work easier to maintain and explain.

\section*{Limitations}
\addcontentsline{toc}{section}{Limitations}
The main limitation of the system is the board itself. The Zybo Legacy is a very useful
platform for learning and prototyping, but the Zynq Z-7010 provides only a limited margin
for implementing CNNs. The final model had to be small and tightly constrained in terms
of resources. This directly affects the accuracy that can be expected and limits the
complexity of the images that can be classified robustly.

Another limitation is that the benchmark focused on inference latency for specific
images. A more complete evaluation would require running a larger test set and obtaining
metrics such as overall accuracy, a confusion matrix, or per-class behavior. In this
work, the priority was to close the full hardware-software workflow and measure the real
acceleration achieved on the board.

Care must also be taken when interpreting the confidence value displayed in the
terminal. The IP core returns logits, not probabilities. The confidence value is
computed in the application using an approximate softmax in order to make the results
easier to read. Therefore, the predicted class and the logits are more relevant than the
exact confidence percentage.

\section*{Future work}
\addcontentsline{toc}{section}{Future work}
A clear line of future work would be to improve the model while preserving the resource
constraints. Techniques such as weight pruning, quantization-aware training,
depthwise separable architectures, or a more systematic search for layer-specific word
lengths could be explored. The goal would be to increase accuracy without exceeding the
LUT and BRAM limits of the board
\citep{han2015deep,howard2017mobilenetsefficientconvolutionalneural}.

It would also be interesting to study DSP usage in more detail. In the final design, DSP
utilization is low, while LUT remain the most critical resource. Future work could
focus on enforcing a more controlled mapping of multiplications to DSP blocks, adjusting
data widths to favor that mapping, and comparing different HLS strategies. It will not
always be possible to move logic into DSP blocks, but this aspect deserves a deeper
exploration.

Another natural improvement would be to automate the generation flow. At present, the
process involves training the model, exporting it to hls4ml, generating the IP,
integrating it into Vivado, exporting the XSA, and rebuilding the Vitis platform. Part
of this flow could be unified through scripts that clean older versions, regenerate the
IP, and prepare a final reproducible folder. This would reduce cache- and path-related
errors, which were a real source of problems during the project.

From an evaluation perspective, the next step would be to test many more images instead
of only individual cases. An application capable of reading images from an SD card, or
even from a camera, would make it possible to evaluate the system in a way that more
closely resembles a real use case. It would also be interesting to measure energy
consumption, since one of the reasons for using FPGA in Edge AI is not only latency,
but also energy efficiency \citep{sze2017efficientprocessingdeepneural}.

Finally, the FPGA accelerator could be compared against a more optimized ARM
implementation, for example using NEON or libraries designed for microcontrollers and
embedded processors \citep{lai2018cmsisnnefficientneuralnetwork}. This would make it
possible to better position the result against a software path specifically optimized
for CPU execution. Even so, the work carried out already demonstrates the main idea:
when the computation is adapted to the hardware and properly integrated into the system,
the FPGA can significantly reduce inference latency even on a resource-constrained
board.
